\subsection[\texttt{C++} side]{\proglang{C++} side}\label{sec:binary_bindings}
The \cpp bindings are written using the \pkg{pybind11}
\citep{pybind11} package.
The main purpose of these bindings is to provide \python
interfaces to selected \lgretl estimators and utility functions
in a standardized form. Each \cpp binding has its \python
counterpart (see Section~\ref{sec:python_side} for details).


\subsubsection{Data management}\label{sec:data_management_c++}
\gretlpy expects a dataset provided as a dictionary with
the following entries:
\begin{itemize}
	\item \pythcode{varname}: A list containing the variable names.
	\item \pythcode{obs\_labels}: A list containing the observation labels.
	\item \pythcode{data\_array}: A \pkg{numpy} array (either column-major
		or row-major format) containing the observations for the variables.
	\item \pythcode{periodicity}: An integer indicating the data's periodicity,
		as defined in \cite{gretl_command}. A brief explanation is as follows:
		\begin{itemize}
			\item For time-series data:
				\begin{itemize}
					\item 1 = annual,
					\item 4 = quarterly,
					\item 12 = monthly,
					\item 52 = weekly,
					\item 5, 6, or 7 = daily,
					\item 24 = hourly.
				\end{itemize}
			\item For panel data: The periodicity represents the number
				of rows per data block. This corresponds to:
				\begin{itemize}
					\item The number of cross-sectional units in the case
						of stacked cross-sections, or
					\item The number of time periods in the case of
						stacked time-series data.
				\end{itemize}
			\item For simple cross-sectional data: The periodicity should
				be set to~1.
		\end{itemize}
\end{itemize}
A dictionary containing data for \gretlpy can be created
manually or with the help of one of the utility functions
described in Section~\ref{sec:data_management_python}.

It is important to note that such a dictionary can be
provided in two general and non-overlapping ways:
\begin{itemize}
	\item \textbf{As an individual dataset}: In this case,
		the dataset is passed directly when calling a \gretlpy
		function, such as an estimator or data-handling function.
	\item \textbf{As a global dataset}: This is used whenever
		an individual dataset is not explicitly provided.
\end{itemize}

\subsubsection{Estimators}\label{sec:estimators_c++}
With \cppcode{namespace py = pybind11}:
%\lstinputlisting[style=cppstyle]{../src/gretl4py.h}

For \textit{ARIMA} estimation (specific case):
\begin{lstlisting}[style=cppstyle]
py::dict *arima (const std::string y, const py::tuple &tup,
								const py::kwargs &kwargs);
\end{lstlisting}

For default case:
\begin{lstlisting}[style=cppstyle]
py::dict *ESTIMATOR (const std::string formula,
								const py::kwargs &kwargs);
\end{lstlisting}
where \cppcode{ESTIMATOR} represents one of the supported
estimator functions (see Section~\ref{sec:estimators_python}).
These bindings return a \python dictionary (see
Section~\ref{sec:data_management_c++} for details) and accept
the following arguments:
\begin{enumerate}
	\item \cppcode{const std::string y}: The name of the dependent
		variable in a dataset (see Section~\error{XXX} for details).
	\item \cppcode{const py::tuple \&tup}: A reference to a \python
		tuple containing the SARIMA specification (see Section~\error{XXX} for details).
	\item \cppcode{const std::string formula}: A formula that constitutes
		an equation (see Section~\error{XXX} for details).
	\item \cppcode{const py::kwargs \&kwargs}: A reference to optional
		arguments (see Section~\error{XXX} for details).
\end{enumerate}
